连续调节优雅联合测量的共同 Bloch 半径，需要使其局部正四面体几何始终与完整正交归一测量基兼容。本文给出适用于全部整数 $n\ge5$ 的优雅联合测量显式连续族，以同一相位构造统一奇数与偶数情形。该族经过指定基准测量，保持测量基的正交归一性与固定局部正四面体方向，使共同 Bloch 半径连续覆盖 $[0,2r_E(n)]$，其中 $r_E(n)=\sqrt3/2^{n-1}$ 为基准半径。在固定相位表示中，本文完整分类由两值实相位生成元产生、经过该基准测量并在其邻域内保持共同半径正四面体几何的相位直线所允许的半径规律。由此证明，对每个 $n\ge5$，上界 $2r_E(n)$ 均可达到且在该模型内最优，并确定 $n=5$ 为该模型允许半径增强的最小规模。对每个偶数 $n\ge8$，本文进一步构造一个二维测量族，其两个坐标截面分别逐列恢复新增强族与已知余弦族。低维结果进一步揭示相位路径形状的作用：在同一相位表示下，当 $n=4$ 时，即使允许任意实相位生成元，所有经过基准测量并在其邻域内保持上述几何条件的相位直线仍不能增强半径，而一条显式弯曲路径可达到 $(1+1/\sqrt2)r_E(4)$。
